\section{Primary geometry at a conductor boundary}
\label{sec:loewy-boundary}
The distinction between a presentation of the full failure scheme and
its primary geometry is essential along the subalgebra locus. We now
resolve that geometry for every complex local algebra with Hilbert
function $(1,2,2)$. The proof begins with a factorization valid for all
algebras whose augmentation ideal has cube zero. In the classified
case the quadratic relation determines the embedded associated strata,
including a stratum supported entirely at the extreme-corank boundary.
No restriction to the quadratic-generation open is made.

\subsection{The quadratic-layer factorization}
Let $B$ be a finite local $\C$-algebra, with maximal ideal
$\mathfrak m$ satisfying $\mathfrak m^3=0$. Put
$V=\mathfrak m/\mathfrak m^2$, $S_2=\mathfrak m^2$, and write
$e=\dim V$, $q=\dim S_2$. Multiplication induces a surjection
\[
 \gamma:\Sym^2V\longrightarrow S_2.
\]
Choose a splitting $\mathfrak m=V\oplus S_2$. The unital parameter
space of $(e+1)$-planes is $X=\Gr(e,\mathfrak m)$: its tautological
plane is $W=\mathcal O\cdot1\oplus K$. On its frame bundle write
\begin{equation}\label{eq:loewy-frame}
 \begin{pmatrix}M\\H\end{pmatrix}:\mathcal O^e\longrightarrow
             V\oplus S_2,
 \qquad M\in\operatorname{Mat}_{e\times e},\quad
 H\in\operatorname{Mat}_{q\times e}.
\end{equation}
The stacked matrix has rank $e$. Its entries, not just the entries
of $M$, are the frame coordinates.

\begin{theorem}[Quadratic-layer factorization]
\label{thm:loewy-factorization}
For every degree $m\ge2$, the full zeroth-Fitting ideal of
$\Sym^m W\to B$ pulls back to
\begin{equation}\label{eq:loewy-factorization}
 \mathcal I_D=(\det M)\,I_q\bigl(\gamma\Sym^2M\bigr).
\end{equation}
The identity holds on the whole frame bundle and after arbitrary
base change, including substitution of coefficients in a nonreduced
ring. In particular it also gives equations at $M=0$ whenever that
stratum meets the frame bundle. The variables $H$ are unrestricted
except for the frame condition; the formula is not a transverse slice.
\end{theorem}
\begin{proof}
Products of three elements of $K$ vanish. Because $1\in W$, the
multiplication image in every degree $m\ge2$ is
$\mathcal O\cdot1+K+K^2$. After deleting zero and repeated columns,
a presentation of its cokernel is given by
\[
 \begin{pmatrix}
 1&0&0\\
 0&M&0\\
 0&H&\gamma\Sym^2M
 \end{pmatrix}.
\]
Every nonzero maximal minor must use the first column and all $e$
columns of the middle block: no other column has an entry in the
$V$ rows. Expansion in the first $e+1$ rows leaves exactly
$\det M$ times a $q$ by $q$ minor of the last block. These are all
the maximal minors. This is a polynomial identity in the universal
coefficient ring. It proves the assertion before taking a radical,
localizing at a tangent minor, or passing to residue fields.
\end{proof}

The same calculation applies to an augmented locally free algebra
$\mathcal O\oplus V\oplus S_2$ over a base scheme, with locally free
layers, $S_2(V\oplus S_2)=0$, and multiplication $\gamma$. Formula
\eqref{eq:loewy-factorization} therefore describes families of quadratic
relations as well as families of generating planes.

\subsection{The complete class of Hilbert function $(1,2,2)$}
Every local algebra in this class has $\mathfrak m^3=0$. The kernel
of $\gamma:\Sym^2V\to S_2$ is a line of binary quadratic forms.
Over $\C$ its generator is either a square or a product of two
independent linear forms. Thus the complete list of isomorphism types is
\begin{equation}\label{eq:loewy-two-types}
 \begin{split}
 B_{\mathrm{sq}}&=\C[x,y]/\bigl((x,y)^3,y^2\bigr),\\
 B_{\mathrm{tf}}&=\C[x,y]/\bigl((x,y)^3,xy\bigr).
 \end{split}
\end{equation}
Indeed a linear splitting of $\mathfrak m\to V$ identifies $B$ with
$\C\oplus V\oplus S_2$, with multiplication determined entirely by
$\gamma$; there are no further extension coefficients when
$\mathfrak m^3=0$.

In both cases $X=\Gr(2,\mathfrak m)$ parametrizes unital
three-planes. On its frame bundle put
\[
 M=\begin{pmatrix}a&c\\b&d\end{pmatrix},\qquad
 \delta=ad-bc,\qquad
 R=\C[a,b,c,d],\quad P=(a,c),\quad Q=(b,d),\quad
 \mathfrak n=(a,b,c,d).
\]
The four entries of $H$ are independent additional variables.
All the following ideals are extended to $R[H]$ and restricted to
the open where the stacked matrix \eqref{eq:loewy-frame} has rank two.
A component absent from an individual affine chart is simply omitted
there; the decomposition is irredundant on the whole frame bundle.

\begin{theorem}[Complete primary decomposition for Hilbert function $(1,2,2)$]
\label{thm:loewy-primary}
For every $m\ge2$ the full multiplication failure scheme has the
following primary decompositions.

For $B_{\mathrm{sq}}$,
\begin{equation}\label{eq:loewy-square-primary}
 I_{\mathrm{sq}}=\delta^2P^2=(\delta^2)\cap P^4,
 \qquad
 \operatorname{Ass}(R/I_{\mathrm{sq}})=\{(\delta),P\}.
\end{equation}
The prime $P$ is embedded. More generally, for every $j\ge1$,
\begin{equation}\label{eq:loewy-square-powers}
 I_{\mathrm{sq}}^j=(\delta^{2j})\cap P^{4j}.
\end{equation}

For $B_{\mathrm{tf}}$, set
\begin{equation}\label{eq:loewy-mixed-ideals}
 \begin{split}
 J&=(ab,ad+bc,cd),\\
 Q_0&=J+P^2+Q^2,\\
 Q_*&=\delta^2Q_0+\mathfrak n^7.
 \end{split}
\end{equation}
Then $Q_0$ and $Q_*$ are $\mathfrak n$-primary and
\begin{equation}\label{eq:loewy-two-factor-primary}
 \begin{split}
 I_{\mathrm{tf}}&=\delta^2J
       =(\delta^2)\cap P^3\cap Q^3\cap Q_*,\\
 \operatorname{Ass}(R/I_{\mathrm{tf}})
       &=\{(\delta),P,Q,\mathfrak n\}.
 \end{split}
\end{equation}
The last three primes are embedded. The displayed $Q_*$ is a
specified primary representative, not a canonical choice of an
embedded primary component.

These identities describe the entire schemes on $X$, including the
extreme-corank locus $T$. Intrinsically, an embedded surface is indexed
by a linear factor $L\subset V$ of the quadratic relation: its support is
\begin{equation}\label{eq:loewy-factor-surface}
 Y_L=\{K\subset\mathfrak m:K\longrightarrow V/L\text{ is zero}\}
       =\Gr(2,S_2+L)\simeq\Pj^2.
\end{equation}
There is one such surface for a square and two for distinct factors.
In the two-factor case there is also the embedded point $K=S_2$.
There are no further associated strata.

The reduced support in both cases is the same integral normal divisor
\[
 \Delta=\{\rank(K\to V)\le1\}\subset X.
\]
Its singular locus is the point $K=S_2$. The generic multiplicity of
$D$ along $\Delta$ is two, whereas the nilradical of $\mathcal O_D$
has exact nilpotency index four. Neither full failure scheme is Cartier.
\end{theorem}
\begin{proof}
Choose bases in the quadratic layer, and rescale the middle symmetric
basis vector where necessary. The two quadratic product matrices are
\[
 L_{\mathrm{sq}}=
 \begin{pmatrix}a^2&2ac&c^2\\ab&ad+bc&cd\end{pmatrix},\qquad
 L_{\mathrm{tf}}=
 \begin{pmatrix}a^2&2ac&c^2\\b^2&2bd&d^2\end{pmatrix}.
\]
For the first matrix the target basis may be taken to be $x^2,2xy$;
for the second it is $x^2,y^2$. Their two-minor ideals are, respectively,
$\delta P^2$ and $\delta J$. Theorem~\ref{thm:loewy-factorization}
gives the asserted full failure ideals. Notice the extra factor
$\delta$ supplied by the linear layer; dropping that layer would
produce a different primary problem.

For the square case, the $P$-adic associated graded ring of $R$ is
$\C[b,d][a,c]$, a domain. The $P$-adic order of $\delta$ is one.
Consequently
\[
 (P^t:\delta^u)=P^{\max\{t-u,0\}}\qquad(t,u\ge0).
\]
It follows that $(\delta^2)\cap P^4=\delta^2P^2$, and the same
calculation with $2j,4j$ proves \eqref{eq:loewy-square-powers}.
The polynomial $\delta$ is prime, and powers of $P$ are primary.
The two components are irredundant: $\delta^2\notin P^4$ and
$a^4\notin(\delta^2)$. This proves the associated-prime assertion.

For the second case we first prove
\begin{equation}\label{eq:loewy-j-decomposition}
 J=P\cap Q\cap Q_0.
\end{equation}
Since $P\cap Q=PQ$, its quotient by $J$ is spanned by the class of
$\delta$: the relations set $ab=cd=0$ and $ad=-bc$. Every element
of $\mathfrak n$ annihilates this class, and the class is nonzero
by linear independence in degree two. In $R/Q_0$ a basis is
\[
 1,\ a,\ b,\ c,\ d,\ ad,
\]
with $bc=-ad$ and all other quadratic monomials zero. Thus
$\mathfrak n^3\subset Q_0$, $Q_0$ is $\mathfrak n$-primary, and
$\delta=2ad$ remains nonzero in this quotient. The natural map
$PQ/J\to R/Q_0$ is therefore injective, proving
\eqref{eq:loewy-j-decomposition}. It is irredundant; the embedded
prime $\mathfrak n$ is also witnessed directly by
$\operatorname{ann}_{R/J}(\bar\delta)=\mathfrak n$.

The same associated-graded argument as before gives
\[
 (P^3:\delta^2)=P,\qquad (Q^3:\delta^2)=Q.
\]
We claim that $(Q_*:\delta^2)=Q_0$. For a homogeneous polynomial
of degree $j$, membership of its product with $\delta^2$ in
$\mathfrak n^7$ is equivalent to $j\ge3$. More generally the
homogeneous decomposition, and injectivity of multiplication by
$\delta^2$ in $R$, give
\[
 (\delta^2Q_0+\mathfrak n^7:\delta^2)
       =Q_0+\mathfrak n^3=Q_0.
\]
Intersecting with $(\delta^2)$ and using
\eqref{eq:loewy-j-decomposition} proves
\eqref{eq:loewy-two-factor-primary}. Since $Q_*\subset\mathfrak n$
and contains $\mathfrak n^7$, it is $\mathfrak n$-primary.
For completeness multiplication by $\delta^2$ gives an injection
$R/J\hookrightarrow R/I_{\mathrm{tf}}$. Hence its three associated
primes $P,Q,\mathfrak n$ remain associated. The prime $(\delta)$
is the unique minimal prime of $I_{\mathrm{tf}}$. This proves
irredundance and excludes additional associated primes.

All these decompositions survive polynomial extension by $H$ and
localization, and every associated stratum meets the frame open
by taking $H$ to be the identity matrix. The global ideals,
irredundancy, and associated points are established formally in
Corollary~\ref{cor:loewy-primary-descent}, using
Lemma~\ref{lem:frame-primary-descent}. That corollary uses only
the affine primary equalities just proved; it does not assume the
global conclusion of the present theorem. In particular no
specialization assertion or unexpanded regular-fibre argument is
needed here.

The reduced equation on the frame cover is $\delta=0$. It is an
integral hypersurface whose singular locus is $M=0$, of codimension
three within that hypersurface. A hypersurface is Cohen--Macaulay;
regularity in codimension one and Serre's criterion give normality.
Integrality on $X$ also follows from the surjection
$\Gr_{\Pj(V)}(2,S_2+\mathcal L)\to\Delta$ from an integral scheme.
At a general point of $\Delta$ the residual ideal $P^2$ or $J$ is a
unit ideal, so the generic multiplicity is two.

In either case write $I=\delta^2J'$, where $J'=P^2$ or $J$.
One has $\delta\notin J'$ and $\delta^2\in J'$; in the latter
case use
$\delta^2=(ad+bc)^2-4ab\,cd$.
Therefore $\delta^4\in I$ and $\delta^3\notin I$.
Since $\sqrt I=(\delta)$, the nilradical has exact index four.
Nonvanishing is retained on the chart where $H$ is invertible.
Finally a nonzero principal ideal in the regular scheme $X$ cannot
have these embedded associated primes. Thus $D$ is not Cartier.
\end{proof}

\subsection{Conductor equations on the embedded strata}
The preceding theorem computes local rings across $T$, not just the
fibres of the incidence over $T$. We now locate the conductor strata
inside these primary supports. This also specifies which embedded
components lie entirely in the extreme-corank locus.

\begin{theorem}[Conductor position of the primary strata]
\label{thm:loewy-conductor-boundary}
On the exact rank-one locus of $K\to V$, write
$L=\operatorname{im}(K\to V)$ and $H_0=K\cap S_2$.
Both are line bundles. The subalgebra locus $T$ there is precisely
the zero scheme of
\begin{equation}\label{eq:loewy-subalgebra-equation}
 \gamma|_{\Sym^2L}:\Sym^2L\longrightarrow S_2/H_0.
\end{equation}
On that scheme the conductor is the actual kernel
\begin{equation}\label{eq:loewy-conductor-kernel}
 \operatorname{cond}_B(W)=
 \ker\bigl(K\longrightarrow
         \operatorname{Hom}(V,S_2/H_0)\bigr),
 \quad k\longmapsto\bigl(v\longmapsto\gamma(\bar k,v)\bigr).
\end{equation}
If the induced map $L\to\operatorname{Hom}(V,S_2/H_0)$ has exact
rank zero, the conductor is $K$ and the conductor action has rank
one. If it has exact rank one, the conductor is $H_0$ and the action
has rank two. These kernel statements commute with base change on
the specified scheme-theoretic rank strata.

For $B_{\mathrm{sq}}$ the embedded surface $Y_{\langle y\rangle}$
is contained in $T$. On its rank-one part the action rank is one
exactly when $H_0=\langle xy\rangle$, and is two otherwise.
For $B_{\mathrm{tf}}$ the two embedded surfaces are
$Y_{\langle x\rangle}$ and $Y_{\langle y\rangle}$.
Their rank-one intersections with $T$ are, respectively,
$H_0=\langle x^2\rangle$ and $H_0=\langle y^2\rangle$;
the conductor action has rank one on both intersections.
Elsewhere on the rank-one part of $T$ it has rank two.

At $K=S_2$, in both types, $W=\C\oplus S_2$ has conductor $S_2$
and quotient $B/S_2=\C\oplus V$ with square-zero augmentation ideal.
The incidence fibre there is $\Pj(V)$, whereas the action-rank-two
fibres above are spectra of the dual numbers. The extra embedded
point at $K=S_2$ occurs precisely in the two-factor type.
Thus numerical action rank alone does not determine associated primes:
the quadratic relation in the conductor data is indispensable.
\end{theorem}
\begin{proof}
Since $S_2\mathfrak m=0$, one has
$W^2=\mathcal O\cdot1+K+\gamma(L^2)$.
The quotient $S_2/H_0$ is a line bundle, so closure under multiplication
is exactly \eqref{eq:loewy-subalgebra-equation}, including its
nonreduced zero-scheme structure.

An element $\lambda\cdot1+k\in W$ in the conductor must have
$\lambda=0$: multiplication by $S_2$ and projection to the free
line $S_2/H_0$ is scalar multiplication by $\lambda$.
For $k\in K$, multiplication by $S_2$ is zero, and multiplication
by a lift of $v\in V$ is $\gamma(\bar k,v)$.
It belongs to $W$ exactly when it belongs to $H_0$.
This proves \eqref{eq:loewy-conductor-kernel} over every coefficient
ring on the indicated rank locus. Its restriction to the exact-rank
strata is a kernel of constant rank, so the sequences split locally
and pull back. The conductor has rank two or one accordingly;
as $\rank W=3$, the action ranks are one or two.

In the square type, $\gamma(y^2)=0$ and
$\gamma(y,V)=\langle xy\rangle$. This proves the assertion about
the entire embedded surface and its action-rank jump. If $L\ne
\langle y\rangle$, then $\gamma(L,V)=S_2$, so action rank one
cannot occur on its rank-one subalgebra locus.
In the two-factor type, multiplication by $x$ has image
$\langle x^2\rangle$ and multiplication by $y$ has image
$\langle y^2\rangle$. For any other line both quadratic basis
vectors occur and the image is $S_2$. These observations prove all
the rank-one assertions. They are equations in the universal line
and quotient bundles, rather than an appeal to continuity of ranks.

At $K=S_2$ the conductor computation is immediate.
The action-rank-one incidence is the scheme of length-two quotients
of $\C\oplus V$, hence the projective line: its binary cubic
square obstruction vanishes identically. On the action-rank-two
locus the rank-two conductor algebra is local with square-zero
augmentation line, since every quadratic product in $K$ lies in
$H_0$. Proposition~\ref{prop:codim-two-incidence-fibres} therefore
gives the spectrum of the dual numbers on geometric fibres.
The assertion about the embedded point follows from
Theorem~\ref{thm:loewy-primary}, not from the incidence fibre list.
\end{proof}

On the rank-one, action-rank-one loci just described, the residual
rank-three algebra is curvilinear: the surviving tangent vector has
nonzero square. Thus its length-two quotient scheme has length three.
The change from these cubic fibres, or the quadratic fibres, to the
projective line at $K=S_2$ is accompanied by the local equations
\eqref{eq:loewy-square-primary} and
\eqref{eq:loewy-two-factor-primary}. Those equations, rather than
the dimensions of the fibres, determine the nilpotent boundary.

\begin{proposition}[Collision of the quadratic factors]
\label{prop:loewy-collision}
The algebras
\[
 \mathcal B_\tau=\C[\tau,x,y]/\bigl((x,y)^3,y^2-\tau xy\bigr)
\]
form a finite free rank-five family over $\C[\tau]$. On the full
unital three-plane frame space its failure ideal in every degree
$m\ge2$ is
\begin{equation}\label{eq:loewy-collision}
 \delta^2\bigl(a(a+\tau b),\
       2ac+\tau(ad+bc),\ c(c+\tau d)\bigr).
\end{equation}
Its fibre at $\tau=0$ is the square-type scheme, and every fibre
at $\tau\ne0$ is the two-factor-type scheme. Thus the primary
strata in Theorem~\ref{thm:loewy-primary} occur in one family of
full Fitting schemes with fixed algebra length and Hilbert function.
No flatness of that failure family, or compatibility of a chosen
primary decomposition with specialization, is asserted.
\end{proposition}
\begin{proof}
The elements $1,x,y,x^2,xy$ are a free basis over $\C[\tau]$;
the displayed relation is the only elimination in degree two and
all terms of degree at least three vanish. The quadratic matrix,
in the product basis without rescaling the middle column, is
\[
 \begin{pmatrix}
 a^2&ac&c^2\\
 2ab+\tau b^2&ad+bc+\tau bd&2cd+\tau d^2
 \end{pmatrix}.
\]
Its minors are $\delta$ times the three displayed generators.
Theorem~\ref{thm:loewy-factorization} proves the ideal identity over
$\C[\tau]$ and after specialization. The quadratic relation is
$y(y-\tau x)$, proving the assertions about the fibres.
\end{proof}

\begin{corollary}[Projective realization of the complete boundary]
\label{cor:loewy-global}
Let $Z\subset\Pj^2$ have homogeneous ideal
$((X,Y)^3,Y^2)$ or $((X,Y)^3,XY)$ in $\C[T,X,Y]$.
For every $n\ge5$, the inverse-image linear series restricting to
all generating three-planes on $Z$ have, in each symmetric degree
$m\ge2$, exactly the primary schemes of
Theorem~\ref{thm:loewy-primary}, on their smooth unit covers.
In particular the embedded conductor-boundary components are
components of actual global section-multiplication failure.
\end{corollary}
\begin{proof}
The ideals are saturated because $T$ is a nonzerodivisor on the
quotient. The affine algebras have top graded degree two, so their
homogeneous ideals have regularity three. Apply
Theorem~\ref{thm:regularity-conductor}, using $n\ge2\cdot3-1$,
and trivialize $\mathcal O_Z(n)$ by $T^n$.
The cokernel isomorphism transfers the full Fitting ideal. Finally
Remark~\ref{rem:codim-two-unit-cover} transfers the equations from
unital planes to a smooth cover of the unrestricted generating
Grassmannian. No boundary is removed in either step.
\end{proof}
